\documentclass[11pt,a4paper]{article}
\usepackage[spanish]{babel}
\usepackage[utf8]{inputenc}
\usepackage[T1]{fontenc}
\usepackage{amsmath}
\usepackage{amssymb}
\usepackage{graphicx}
\usepackage{booktabs}
\usepackage{float}
\usepackage{hyperref}
\usepackage[numbers,sort&compress]{natbib}
\usepackage{geometry}
\geometry{margin=2.5cm}

% Notación del proyecto
\newcommand{\Pscint}{P_{\mathrm{scint}}}
\newcommand{\kioq}{k_{\mathrm{ioq}}}
\newcommand{\Cmaterial}{C_{\mathrm{material}}}
\newcommand{\kB}{k_B}
\newcommand{\Lideal}{L_{\mathrm{ideal}}}
\newcommand{\Lquenched}{L_{\mathrm{quenched}}}
\newcommand{\Dscint}{D_{\mathrm{scint}}}
\newcommand{\Dscintw}{D_{\mathrm{scint}(w)}}
\newcommand{\kQclin}{k_{Q_{\mathrm{clin}},Q_{\mathrm{msr}}}^{f_{\mathrm{clin}},f_{\mathrm{msr}}}}

\title{Quenching de ionización en detectores de centelleo plástico:\\
Estado del arte y posicionamiento de MDSIM1}
\author{Paul Martinez}
\date{Mayo 2026}

\begin{document}
\maketitle
\tableofcontents
\newpage

% ─────────────────────────────────────────
\section{El detector ideal y su defecto oculto}
% Introducción: PSD, agua-equivalencia, campos pequeños.
% Quenching como problema central. Gancho narrativo.
% Refs: beddar1992, palmans2018, gingras2025

% CONTENIDO AQUÍ



Un detector de centelleo plástico (PSD, por \textit{plastic scintillation detector}) consiste en un cilindro de material orgánico centelleador, típicamente de $\sim$1~mm de diámetro y 1~mm de longitud, acoplado a una fibra óptica que guía la señal luminosa hasta un sistema de lectura optoelectrónico~\cite{beddar1992}. Cuando la radiación ionizante deposita energía en el volumen sensible, las moléculas del centelleador se excitan y emiten fotones de fluorescencia; la intensidad de esa emisión es, en principio, proporcional a la dosis absorbida. Esta premisa ha sostenido más de tres décadas de desarrollo en dosimetría clínica.

Las ventajas del PSD como dosímetro en radioterapia son difíciles de igualar con una sola tecnología alternativa. Su composición --- fundamentalmente carbono e hidrógeno con $Z_{\mathrm{eff}} \approx 5{.}7$ --- le confiere una agua-equivalencia superior a la de diodos de silicio, detectores de diamante o cámaras de ionización en términos de cocientes de poder de frenado másico y coeficientes de absorción másicos de energía~\cite{beddar1992}. Su volumen sensible sub-milimétrico minimiza el promediado volumétrico de dosis, un artefacto que distorsiona las mediciones cuando el gradiente de dosis es pronunciado. No requiere corrección por tasa de dosis, presenta excelente linealidad en el rango terapéutico y ofrece lectura en tiempo real~\cite{bauer2026}. Estas propiedades lo convierten en un candidato natural para tres de los escenarios más exigentes de la dosimetría moderna: la dosimetría de campos pequeños, la dosimetría en superficie y la dosimetría \textit{in vivo}.

El contexto clínico de los campos pequeños justifica por sí solo la búsqueda de un detector con estas características. Técnicas como la radiocirugía estereotáctica, la radioterapia de intensidad modulada y los tratamientos con colimador de microláminas emplean campos con dimensiones del orden de $0{.}5 \times 0{.}5$~cm$^{2}$ a $4 \times 4$~cm$^{2}$. En estos campos se pierde el equilibrio lateral de partículas cargadas, la fuente de radiación puede quedar parcialmente ocluida y el gradiente de dosis en los bordes es extremadamente abrupto~\cite{palmans2018}. Los detectores convencionales --- cámaras de ionización de pequeño volumen, diodos, detectores de diamante --- sufren en este régimen perturbaciones que obligan a aplicar factores de corrección específicos para cada combinación de detector, tamaño de campo y calidad de haz. El formalismo del código de práctica IAEA TRS-483 establece que el factor de salida de campo medido debe corregirse mediante el factor $k^{f_{\mathrm{clin}},f_{\mathrm{msr}}}_{Q_{\mathrm{clin}},Q_{\mathrm{msr}}}$, que puede desviarse significativamente de la unidad para detectores de estado sólido en los campos más pequeños~\cite{palmans2018}. Por el contrario, el TRS-483 prescribía que el factor de corrección de los PSD podía asumirse igual a la unidad, dado su alto grado de agua-equivalencia.

Esta presunción de idealidad se mantuvo como consenso durante más de una década. Y sin embargo, contenía un defecto oculto.

Los centelleadores plásticos están sujetos al fenómeno de \textit{ionization quenching}: una disminución no lineal de la producción de luz por unidad de dosis absorbida cuando la deposición de energía es causada por partículas cargadas con alta densidad de ionización~\cite{beddar1992,santurio2019}. En haces de partículas pesadas o radiación de kilovoltaje, el efecto era conocido y corregido. En haces de fotones de megavoltaje, sin embargo, se asumía implícitamente que no existía~\cite{palmans2018}. El trabajo de Valdes Santurio y Andersen~\cite{santurio2019} demostró que esta suposición era incorrecta: incluso en un haz de 6~MV, entre el 33\% y el 44\% de la dosis absorbida en el centelleador es depositada por electrones secundarios con energía inferior a 125~keV, un umbral por debajo del cual el \textit{quenching} se vuelve apreciable. Dado que la fracción de estos electrones de baja energía cambia con el tamaño de campo y con la calidad del haz, el \textit{quenching} introduce una dependencia sistemática que permanecía bajo la sombra del presupuesto de incertidumbre de las mediciones clínicas.

El hallazgo reciente de Gingras \textit{et al.}~\cite{gingras2025} ha añadido un giro inesperado a esta historia. Trabajando con el detector PRB-0002 de la plataforma Hyperscint RP-200, los autores realizaron la primera separación completa de los factores de perturbación de un PSD comercial para campos entre $0{.}6 \times 0{.}6$ y $30 \times 30$~cm$^{2}$ en un haz de 6~MV de un acelerador Varian TrueBeam. Los resultados muestran que el factor de perturbación por composición del centelleador, $P_{\mathrm{scint}}$, y el factor de corrección por \textit{quenching} de ionización, $k_{\mathrm{ioq}}$, se compensan mutuamente, produciendo un factor de corrección total $k^{f_{\mathrm{clin}},f_{\mathrm{msr}}}_{Q_{\mathrm{clin}},Q_{\mathrm{msr}}}$ entre 0.999 y 1.002 en todo el rango de tamaños de campo. El detector se comporta como si fuera ideal --- pero no porque carezca de defecto, sino porque dos imperfecciones de signo opuesto se cancelan.

Esta compensación accidental plantea una pregunta fundamental: ¿es la cancelación entre $P_{\mathrm{scint}}$ y $k_{\mathrm{ioq}}$ una propiedad universal de los centelleadores plásticos, o depende de la composición específica del material, de la geometría del detector y de la calidad del haz? Si depende del material, entonces el \textit{quenching} no es un efecto despreciable que se puede ignorar, sino un parámetro de diseño que debe controlarse. Y si es un parámetro de diseño, entonces la simulación Monte Carlo se convierte en la herramienta indispensable para explorar el espacio de materiales antes de construir un prototipo.

% ─────────────────────────────────────────
\section{¿Qué es el quenching? Del centelleo a la saturación}
% Mecanismo físico sin ecuaciones. Para audiencia clínica.
% Centelleador: radiación → excitación → fotón.
% Qué interrumpe ese camino. Umbral 125 keV.
% Refs: beddar1992, birks1951, santurio2019

% CONTENIDO AQUÍ


Para comprender el \textit{quenching} de ionización es necesario primero entender cómo un centelleador plástico produce luz. El proceso puede describirse en tres etapas sucesivas. En la primera, una partícula cargada --- típicamente un electrón secundario puesto en movimiento por un fotón del haz de tratamiento --- atraviesa el material orgánico y deposita energía a lo largo de su trayectoria, ionizando y excitando las moléculas del plástico base (poliestireno o polivinilueno). En la segunda etapa, esa energía de excitación se transfiere de las moléculas de la matriz a las moléculas de un dopante fluorescente disperso en el material; el dopante absorbe la energía y la re-emite como fotones de fluorescencia con una longitud de onda desplazada respecto a la absorción, lo que permite que la luz viaje distancias apreciables sin ser reabsorbida. En la tercera etapa, los fotones de centelleo se propagan a través de una fibra óptica acoplada al volumen sensible hasta alcanzar un fotodetector que convierte la señal luminosa en una señal eléctrica~\cite{beddar1992}.

En condiciones ideales, la intensidad de la señal luminosa es directamente proporcional a la dosis absorbida en el centelleador. Esta linealidad es la base de toda la dosimetría con PSD: si la eficiencia de centelleo --- definida como la fracción de energía depositada que se convierte en luz de fluorescencia --- es constante, entonces medir luz equivale a medir dosis~\cite{beddar1992}. Sin embargo, la eficiencia de centelleo no es constante en todas las condiciones.

El \textit{ionization quenching} es el fenómeno por el cual la producción de luz por unidad de dosis absorbida disminuye cuando la densidad de ionización a lo largo de la trayectoria de la partícula cargada es elevada~\cite{birks1951}. El mecanismo físico subyacente puede entenderse sin recurrir a ecuaciones. Cuando un electrón de alta energía (por ejemplo, 1~MeV) atraviesa el centelleador, deposita su energía de forma dispersa: las moléculas excitadas están suficientemente separadas unas de otras como para que cada una complete su ciclo de fluorescencia de manera independiente, y la luz producida es proporcional a la energía depositada. Pero cuando un electrón de baja energía (por ejemplo, 10~keV) recorre el mismo material, su poder de frenado es mucho mayor: deposita la misma energía en una distancia mucho más corta, creando una densidad local de moléculas excitadas e ionizadas muy elevada. En esta región de alta densidad de excitación, las moléculas dañadas ocupan estados altamente excitados o ionizados y se desexcitan rápidamente por vías no radiativas --- es decir, disipan su energía como calor o vibraciones de la red molecular en lugar de emitir fotones~\cite{boivin2016}. Estos procesos de desexcitación no radiativa compiten directamente con la fluorescencia, y el resultado neto es una pérdida de señal luminosa: el centelleador absorbe la misma dosis, pero produce menos luz.

La consecuencia práctica es que la relación entre luz y dosis deja de ser lineal cuando una fracción significativa de la energía es depositada por partículas con alto poder de frenado. Birks~\cite{birks1951} fue el primero en formalizar esta observación en 1951, introduciendo un modelo semiempírico que describe la producción de luz específica (fotones por unidad de longitud de trayectoria) como una función de la pérdida de energía específica de la partícula. En el modelo de Birks, la relación entre luz y energía depositada es lineal a bajos valores de $\mathrm{d}E/\mathrm{d}x$ (electrones de alta energía, régimen sin \textit{quenching}) y se satura progresivamente a medida que $\mathrm{d}E/\mathrm{d}x$ aumenta (electrones de baja energía, protones, iones pesados). La transición entre ambos regímenes está gobernada por una constante del material, $kB$, conocida como el coeficiente de Birks. El valor de $kB$ depende de la composición del centelleador, no del tipo de partícula incidente, y caracteriza la susceptibilidad del material al \textit{quenching}~\cite{boivin2016}.

El umbral práctico por debajo del cual el \textit{quenching} se vuelve relevante para electrones en centelleadores orgánicos es de aproximadamente 125~keV de energía cinética. Por encima de este valor, la respuesta del centelleador a electrones es esencialmente lineal con la energía depositada, y las extrapolaciones lineales desde altas energías pasan a pocos keV del origen~\cite{beddar1992}. Por debajo de 125~keV, la eficiencia de centelleo cae de forma apreciable. Beddar \textit{et al.}~\cite{beddar1992} argumentaron en 1992 que esta caída tenía muy poco efecto en el contexto de los haces de megavoltaje utilizados en radioterapia, dado que el espectro de electrones secundarios en esos haces está dominado por partículas de alta energía. Esta afirmación fue razonable en su momento y se convirtió en la justificación implícita para ignorar el \textit{quenching} en dosimetría con PSD durante más de dos décadas.

Lo que no se había cuantificado hasta el trabajo de Valdes Santurio y Andersen~\cite{santurio2019} era la fracción real de dosis depositada por electrones de baja energía en un haz clínico de megavoltaje. Sus simulaciones Monte Carlo mostraron que, en un haz de 6~MV con un campo de $10 \times 10$~cm$^{2}$, entre el 33\% y el 44\% de la dosis absorbida en el centelleador es atribuible a electrones con energía inferior a 125~keV. Este resultado transformó la narrativa: el \textit{quenching} no era un efecto confinado a los haces de partículas pesadas o a la radiación de kilovoltaje, sino un fenómeno presente --- aunque de magnitud pequeña --- en las condiciones clínicas más comunes de la radioterapia con fotones. Más aún, dado que la fracción de electrones de baja energía depende del tamaño de campo (la radiación dispersa aporta fotones de menor energía que generan electrones secundarios más lentos) y de la calidad del haz, el \textit{quenching} introduce una dependencia sistemática que varía con las condiciones de irradiación.

%En síntesis, el \textit{quenching} de ionización es una interrupción del camino que conecta la energía absorbida con la señal luminosa: un cortocircuito molecular que desvía parte de la energía hacia canales no radiativos antes de que pueda convertirse en fotones detectables. Su magnitud depende de cuán densamente se deposita la energía, y esa densidad depende tanto de la energía de la partícula cargada como de las propiedades del material centelleador. La sección siguiente formaliza esta relación mediante la ley de Birks y sus extensiones.

% ─────────────────────────────────────────
\section{La física del fotón que nunca nació: ley de Birks}
% Ecuación de Birks. Curva dL/dx vs dE/dx.
% Región lineal vs saturación. Valores de kB por material.
% Modelo modificado Santurio 2020.
% Refs: birks1951, santurio2019, santurio2020, ebenau2016

% CONTENIDO AQUÍ


La observación cualitativa --- que la producción de luz se reduce cuando la densidad de ionización es alta --- fue formalizada por Birks en 1951 mediante un modelo semiempírico que sigue siendo la base de la descripción del \textit{quenching} en centelleadores orgánicos~\cite{birks1951}. La ley de Birks relaciona la luz emitida por unidad de longitud de trayectoria, $\mathrm{d}L/\mathrm{d}x$, con la pérdida de energía específica de la partícula cargada, $\mathrm{d}E/\mathrm{d}x$:

\begin{equation}
\frac{\mathrm{d}L}{\mathrm{d}x} = \frac{S \,\dfrac{\mathrm{d}E}{\mathrm{d}x}}{1 + kB \,\dfrac{\mathrm{d}E}{\mathrm{d}x}}
\label{eq:birks}
\end{equation}

\noindent donde $S$ es la eficiencia absoluta de centelleo (fotones por unidad de energía depositada en ausencia de \textit{quenching}) y $kB$ es el coeficiente de Birks, un parámetro empírico del material con unidades de longitud por energía (típicamente cm\,MeV$^{-1}$ o g\,cm$^{-2}$\,MeV$^{-1}$). El modelo trata el \textit{quenching} como un fenómeno unimolecular: la constante $B$ representa la densidad específica de moléculas ionizadas y excitadas a lo largo de la trayectoria de la partícula, y $k$ es la probabilidad de que esas moléculas se desexciten por vías no radiativas. Dado que $k$ y $B$ aparecen siempre como producto, actúan como un único parámetro del material~\cite{boivin2016}.

El comportamiento de la ecuación~\eqref{eq:birks} define dos regímenes claramente diferenciados. Cuando $\mathrm{d}E/\mathrm{d}x$ es bajo (electrones de alta energía, por encima de $\sim$125~keV), el término $kB \cdot \mathrm{d}E/\mathrm{d}x$ en el denominador es mucho menor que la unidad y la ecuación se reduce a $\mathrm{d}L/\mathrm{d}x \approx S \cdot \mathrm{d}E/\mathrm{d}x$: la producción de luz es lineal con la energía depositada, el centelleador responde de forma ideal y $kB$ es irrelevante. Cuando $\mathrm{d}E/\mathrm{d}x$ es alto (electrones de baja energía, protones, iones pesados), el denominador crece y la producción de luz se satura: el centelleador absorbe más energía por unidad de longitud pero no la convierte proporcionalmente en fotones. La curva $\mathrm{d}L/\mathrm{d}x$ \textit{versus} $\mathrm{d}E/\mathrm{d}x$ pasa de ser una recta con pendiente $S$ a curvarse hacia una asíntota horizontal en $S/kB$. Todo fotón que ``debería haber nacido'' según la proporcionalidad lineal, pero que no nace porque la energía se disipa por vías no radiativas, es un fotón perdido por \textit{quenching}.

Un aspecto crucial del modelo de Birks, y fuente de considerable confusión en la literatura, es la determinación experimental de $kB$. En principio, $kB$ es una constante del material centelleador, independiente del tipo de partícula incidente. En la práctica, los valores reportados para un mismo material presentan una dispersión notable. Para centelleadores basados en polivinilueno (PVT), la literatura reporta valores de $kB$ entre 0.002 y 0.021~g\,cm$^{-2}$\,MeV$^{-1}$~\cite{boivin2016}. Para centelleadores basados en poliestireno (PS), el rango publicado es aún más amplio: entre 0.0005 y 0.094~g\,cm$^{-2}$\,MeV$^{-1}$~\cite{boivin2016}. De forma notable, estos dos extremos han sido reportados para el mismo centelleador, la fibra BCF-12~\cite{boivin2016}. Esta dispersión no refleja necesariamente una variabilidad real del material, sino las diferencias en las condiciones experimentales y los métodos de extracción del parámetro: el valor de $kB$ recuperado depende de si se asume que la partícula se detiene completamente dentro del centelleador, del modelo de poder de frenado utilizado, del rango de energías cubierto por el ajuste y de si se consideran o no los electrones secundarios~\cite{santurio2020}.

Esta última observación motivó la extensión más importante del formalismo de Birks para la dosimetría con PSD: el modelo modificado de Valdes Santurio y Andersen~\cite{santurio2020}. La formulación original de la ecuación~\eqref{eq:birks} asume implícitamente que el centelleador es suficientemente grueso como para que tanto el electrón primario como todos sus electrones secundarios se detengan dentro del volumen sensible. En un PSD clínico con dimensión característica de 1~mm, esta asunción no se cumple: una fracción significativa de los electrones secundarios generados dentro del centelleador escapan del volumen sensible sin detenerse, y electrones generados fuera del centelleador pueden entrar y depositar energía. El modelo modificado resuelve esta limitación integrando la ley de Birks explícitamente sobre el espectro completo de electrones (primarios y secundarios) que depositan energía en el volumen sensible, utilizando el poder de frenado electrónico restringido $L_{\Delta}(E)$ con un umbral de $\Delta = 1$~keV:

\begin{equation}
L = \sum_{i=1}^{n} \int_{E_{\min}}^{E_{\max}} \frac{S}{1 + kB \cdot L_{\Delta}(E_i)} \,\mathrm{d}E_i
\label{eq:birks_modified}
\end{equation}

\noindent donde $n$ es el número total de electrones que cruzan, se detienen o se crean en el centelleador, y $L_{\Delta}(E_i)$ es el poder de frenado restringido del $i$-ésimo electrón con energía cinética $E_i$~\cite{santurio2020}. La ventaja fundamental de esta formulación es que no requiere la asunción de equilibrio de partículas cargadas: cada deposición de energía es ponderada individualmente según su densidad de ionización local, lo que permite evaluar el \textit{quenching} en configuraciones donde el equilibrio no se establece, como en los campos pequeños o en las proximidades de interfaces.

Valdes Santurio y Andersen implementaron la ecuación~\eqref{eq:birks_modified} como un objeto \textit{ausgab} dentro del código Monte Carlo EGSnrc y validaron el modelo contra mediciones experimentales del centelleador BCF-60 en haces de rayos X de kilovoltaje (38--158~keV de energía media) en el laboratorio primario italiano ENEA-INMRI~\cite{santurio2020}. El valor de $kB$ obtenido para BCF-60 fue de 0.019~cm\,MeV$^{-1}$, consistente tanto en el régimen de kilovoltaje como en el de megavoltaje~\cite{santurio2019,santurio2020}. Este resultado proporcionó la primera demostración de que un único valor de $kB$, combinado con un modelo de transporte de electrones suficientemente detallado, puede describir el \textit{quenching} en un rango amplio de calidades de haz.

La implementación del modelo de Birks en simulaciones Monte Carlo, sin embargo, introduce una complicación técnica documentada por Dietz-Laursonn~\cite{ebenau2016}. La ecuación~\eqref{eq:birks} es no lineal en $\mathrm{d}E/\mathrm{d}x$: la corrección por \textit{quenching} aplicada a un paso de simulación con un valor promedio de $\mathrm{d}E/\mathrm{d}x$ no es idéntica a la suma de correcciones aplicadas a múltiples sub-pasos con valores fluctuantes de $\mathrm{d}E/\mathrm{d}x$. Como consecuencia matemática, el $kB$ efectivo recuperado de una simulación depende del tamaño de paso del transporte electrónico utilizado por el código Monte Carlo~\cite{ebenau2016}. Un paso grueso promedia las fluctuaciones de $\mathrm{d}E/\mathrm{d}x$ y subestima el \textit{quenching}; un paso fino captura mejor la variabilidad local pero incrementa el costo computacional. Esta dependencia es específica de la implementación: Geant4 aplica la corrección de Birks paso a paso a través de la clase \texttt{G4EmSaturation}, mientras que EGSnrc la aplica mediante un \textit{ausgab object} que puede utilizar el poder de frenado restringido~\cite{santurio2020,ebenau2016}. Los dos enfoques no son equivalentes sin una validación cruzada cuidadosa, y un valor de $kB$ extraído con un código no puede transferirse directamente al otro sin verificar la convergencia del resultado con el control de paso.

%Esta advertencia técnica tiene consecuencias directas para el proyecto MDSIM1: la simulación del $k_{\mathrm{ioq}}$ en Geant4 debe incluir un estudio de convergencia que demuestre que el resultado es estable frente a variaciones del tamaño de paso, antes de que cualquier valor numérico pueda interpretarse como una predicción del \textit{quenching} del detector simulado.

%En síntesis, la ley de Birks proporciona el marco cuantitativo mínimo para describir el \textit{quenching} de ionización: un parámetro del material ($kB$), una ecuación semiempírica y dos regímenes (lineal y saturado). El modelo modificado de Valdes Santurio extiende este marco a detectores sub-milimétricos donde los electrones no se detienen dentro del volumen sensible. Sin embargo, el modelo de Birks --- en cualquiera de sus versiones --- trata el \textit{quenching} como un fenómeno fenomenológico gobernado por un parámetro ajustable, sin especificar el mecanismo molecular que lo causa. La sección siguiente explora un enfoque alternativo, basado en la transferencia resonante de energía de Förster, que intenta dar cuenta del \textit{quenching} desde primeros principios.

% ─────────────────────────────────────────
\section{Más allá de Birks: el mecanismo molecular FRET}
% Ogawa 2018. FRET. R^{-6}. kB no es constante fundamental.
% Kikuta 2021. Implicación para diseño de material.
% Refs: ogawa2018, kikuta2021, laplace2022

% CONTENIDO AQUÍ


El modelo de Birks y sus extensiones --- Chou, Wright, Voltz, Yoshida --- comparten una característica fundamental: describen la reducción de la producción de luz como función de la pérdida de energía específica $\mathrm{d}E/\mathrm{d}x$ de la partícula cargada, sin especificar el mecanismo molecular que causa esa reducción. El coeficiente $kB$ es un parámetro empírico que condensa toda la física microscópica en un solo número ajustable. Esta aproximación ha sido extraordinariamente útil durante siete décadas, pero presenta limitaciones que se han vuelto evidentes a medida que se acumulan datos experimentales para diferentes materiales y tipos de partícula.

La primera limitación es de alcance: el modelo canónico de Birks no describe adecuadamente la producción de luz a bajas energías de retroceso ni para partículas de distinta carga. Von Krosigk \textit{et al.} demostraron que las relaciones de producción de luz para protones y partículas alfa en un centelleador de alquilbenceno lineal no pueden describirse con el mismo valor de $kB$. Yoshida \textit{et al.} reportaron el mismo fenómeno para el centelleador líquido de KamLAND. La evaluación sistemática de Laplace \textit{et al.}~\cite{laplace2022} confirmó que el modelo de Birks falla consistentemente para energías de protones por debajo de $\sim$100~keV en cuatro materiales orgánicos diferentes (EJ-309, EJ-204, EJ-276 y un vidrio orgánico), y que ninguno de los cinco modelos evaluados --- Birks, Chou, Hong, Yoshida y Voltz --- logra describir simultáneamente la producción de luz de protones y de iones de carbono con un conjunto único de parámetros.

La segunda limitación es de universalidad: Pöschl \textit{et al.}~\cite{poeschl2020} realizaron la primera medición modelo-independiente de la función de \textit{quenching} usando un detector de fibra centelleadora segmentado expuesto a protones de 30--100~MeV. Sus resultados para dos materiales, SCSF-78 y BC-408, revelaron que la forma de la función de \textit{quenching} depende fuertemente del material centelleador. Para SCSF-78, el modelo de Chou describe el \textit{quenching} mediante un término cuadrático puro, mientras que para BC-408 lo describe mediante un término lineal puro --- con el término complementario compatible con cero en cada caso. De forma análoga, el modelo de Voltz requiere que la fracción de energía no sujeta a \textit{quenching} ($f$) sea esencialmente cero para SCSF-78, pero cercana al 50\% para BC-408. Estos resultados demuestran que $kB$ no es simplemente un número que varía entre materiales: la \textit{forma funcional} de la dependencia del \textit{quenching} con $\mathrm{d}E/\mathrm{d}x$ es ella misma material-dependiente.

Estas observaciones sugieren que los modelos basados en $\mathrm{d}E/\mathrm{d}x$ como variable única están incompletos. Lo que falta es una descripción que considere explícitamente la distribución espacial de las moléculas excitadas a escala nanométrica, porque es la distancia entre ellas --- no solo la densidad promedio de ionización --- lo que determina si la energía de excitación se emite como un fotón o se disipa por vías no radiativas.

Ogawa \textit{et al.}~\cite{ogawa2018} propusieron en 2018 un modelo radicalmente diferente que aborda exactamente esta cuestión. En lugar de parametrizar el \textit{quenching} como función de $\mathrm{d}E/\mathrm{d}x$, su enfoque se basa en dos elementos: una simulación de estructura de traza que calcula la posición espacial de cada deposición de energía a escala nanométrica, y el efecto de transferencia resonante de energía de Förster (FRET, \textit{Förster Resonance Energy Transfer}) como mecanismo físico del \textit{quenching}.

El mecanismo FRET funciona de la siguiente manera. Cuando la radiación ionizante deposita energía en el centelleador, produce moléculas excitadas --- denominadas ``donoras'' en la terminología de Förster --- que normalmente emitirían fotones de fluorescencia. Sin embargo, si una molécula donora tiene otra molécula excitada (o dañada) suficientemente cerca, puede transferirle su energía de excitación de forma no radiativa mediante un acoplamiento dipolo-dipolo. La probabilidad de esta transferencia depende críticamente de la distancia $r$ entre las dos moléculas, siguiendo la ley:

\begin{equation}
p(r) = \frac{1}{1 + \left(\dfrac{r}{R_f}\right)^{6}}
\label{eq:fret}
\end{equation}

\noindent donde $R_f$ es el radio de Förster, una constante específica del material que caracteriza la distancia a la cual la probabilidad de transferencia es del 50\%~\cite{ogawa2018}. La dependencia con $r^{-6}$ hace que la transferencia sea extremadamente sensible a la distancia: a $r = R_f$ la probabilidad es 0.5, pero a $r = 2R_f$ cae a 0.015. Para los materiales orgánicos utilizados en centelleadores de detección, los radios de Förster reportados se encuentran en el rango de 2 a 6~nm~\cite{ogawa2018}.

La conexión con el \textit{quenching} es directa: cuando una partícula con alto $\mathrm{d}E/\mathrm{d}x$ atraviesa el centelleador, las moléculas excitadas se generan con separaciones del orden de unos pocos nanómetros en el núcleo de la traza, es decir, dentro del radio de Förster. En esas condiciones, una fracción significativa de las moléculas donoras transfiere su energía a vecinas por FRET en lugar de emitir fotones, reduciendo la producción de luz. Para una partícula con bajo $\mathrm{d}E/\mathrm{d}x$, las moléculas excitadas están mucho más separadas, la probabilidad de FRET es despreciable y la luz se produce sin pérdida. El modelo FRET reproduce así el comportamiento fenoménico de Birks, pero lo hace desde la geometría espacial de la deposición de energía, sin recurrir a un parámetro libre de ajuste empírico (el único parámetro es $R_f$, que tiene significado físico medible).

Ogawa \textit{et al.} implementaron este modelo utilizando el código de estructura de traza RITRACKS para calcular las coordenadas $(x,y,z)$ de cada deposición de energía en un centelleador NE-102A (equivalente a EJ-212 y BC-400) irradiado con electrones, protones e iones pesados desde $^{4}$He hasta $^{81}$Br. Asumiendo $R_f = 4$~nm para NE-102A, obtuvieron una concordancia con datos experimentales de producción de luz que se extiende desde electrones de 0.2~MeV hasta iones de bromo de 120~MeV --- un rango que abarca variaciones de más de un orden de magnitud en la producción de luz normalizada~\cite{ogawa2018}. Notablemente, el modelo reproduce una observación que la ley de Birks no puede explicar: para partículas con el mismo LET pero diferente especie iónica, la producción de luz puede diferir, porque la distribución espacial de las excitaciones depende no solo de $\mathrm{d}E/\mathrm{d}x$ sino de la estructura de la traza completa, incluyendo los rayos $\delta$.

La implicación más profunda del modelo FRET para el diseño de detectores es que $kB$ no es una constante fundamental del material centelleador, sino una cantidad emergente que depende de la microestructura del material: la concentración de dopante, la eficiencia de transferencia de energía entre la matriz y el fluoróforo, y la distribución espacial de los centros de centelleo~\cite{ogawa2018}. Cambiar la concentración de dopante, el tipo de polímero base o la arquitectura molecular del centelleador altera $R_f$ y, por tanto, la susceptibilidad al \textit{quenching}. Esta perspectiva transforma el \textit{quenching} de un problema a corregir en un parámetro a diseñar.

La evaluación comparativa de Laplace \textit{et al.}~\cite{laplace2022} proporciona evidencia independiente en esta dirección. Los autores concluyen que los modelos basados en estructura de traza representan el camino más prometedor para superar las limitaciones de los modelos semiempíricos, y que se requiere trabajo adicional para extender estos modelos a iones de carbono y para incorporar la no-linealidad de la respuesta electrónica a bajas energías. El modelo de Voltz --- que introduce una separación explícita entre las componentes rápida y lenta de la centelleación con tratamiento especializado de los rayos $\delta$ --- constituye un puente teórico entre los modelos puramente fenomenológicos y el enfoque de estructura de traza, pero no resuelve la dependencia con la especie iónica.

%Para el proyecto MDSIM1, el modelo FRET no se implementará directamente: su costo computacional (simulación de estructura de traza a escala nanométrica) es incompatible con el volumen de estadística necesario para calcular $k_{\mathrm{ioq}}$ en haces clínicos. Sin embargo, la perspectiva FRET es esencial por tres razones. Primera, justifica por qué el $kB$ de 0.019~cm\,MeV$^{-1}$ medido para BCF-60 no puede transferirse automáticamente al material de Blue Physics sin validación: si la composición del material es diferente, $R_f$ es diferente y $kB$ es diferente. Segunda, fundamenta la hipótesis de que la compensación entre $P_{\mathrm{scint}}$ y $k_{\mathrm{ioq}}$ observada por Gingras~\cite{gingras2025} puede depender del material --- porque $k_{\mathrm{ioq}}$ depende de $kB$, que depende de $R_f$, que depende de la composición molecular. Tercera, conecta el proyecto de Camilo con el de Ian Rossi: el diseño de un nuevo material centelleador para un detector portátil de estado sólido puede beneficiarse de criterios de selección basados en $R_f$ y en la métrica $C_{\mathrm{material}}$ que MDSIM1 propone evaluar.
% ─────────────────────────────────────────
\section{El espectro que nadie miraba: quenching en haces MV}
% Santurio 2019: 33-44\% electrones bajo 125 keV.
% Dependencia con tamaño de campo. Corrección 0.6-2\%.
% Benmakhlouf 2017: mecanismo espectral.
% Refs: santurio2019, santurio2020, benmakhlouf2017, boivin2016

% CONTENIDO AQUÍ


Durante más de dos décadas, la comunidad de dosimetría con centelleadores plásticos operó bajo una asunción implícita: el \textit{quenching} de ionización era un problema de los haces de partículas pesadas y de la radiación de kilovoltaje, pero no de los haces de fotones de megavoltaje utilizados en la práctica clínica. El código de práctica TRS-483~\cite{palmans2018} reflejaba este consenso al prescribir un factor de corrección de salida de campo igual a la unidad para detectores de centelleo plástico en haces de fotones MV, asumiendo que su respuesta era esencialmente agua-equivalente e independiente del tamaño de campo y de la calidad del haz. La asunción tenía una base razonable: Beddar \textit{et al.}~\cite{beddar1992} habían mostrado en 1992 que la respuesta de los centelleadores orgánicos a electrones es lineal por encima de $\sim$125~keV, y el espectro de electrones secundarios en un haz de megavoltaje está dominado por partículas con energías muy superiores a ese umbral. Lo que nadie había cuantificado era la cola de baja energía de ese espectro.

El trabajo de Valdes Santurio y Andersen~\cite{santurio2019} cambió esta perspectiva al realizar la primera cuantificación sistemática de la contribución de los electrones de baja energía a la dosis absorbida en un centelleador plástico irradiado por haces clínicos de fotones MV. Utilizando simulaciones Monte Carlo con EGSnrc y los \textit{phase space files} del acelerador Varian TrueBeam, los autores calcularon la dosis acumulativa como función de la energía de corte de los electrones para un centelleador BCF-60 en distintas configuraciones de campo y calidad de haz. El resultado fue revelador: en un haz de 6~MV con un campo de $10 \times 10$~cm$^{2}$, entre el 33\% y el 44\% de la dosis absorbida en el centelleador proviene de electrones con energía cinética inferior a 100~keV~\cite{santurio2019}. No se trata de una cola despreciable: es más de un tercio de la dosis total.

El origen físico de esta fracción de baja energía es bien conocido en la física de haces, aunque sus consecuencias para el \textit{quenching} no habían sido analizadas. Un haz de fotones de megavoltaje genera electrones secundarios predominantemente por efecto Compton. Estos electrones primarios tienen energías del orden de MeV, pero a medida que se frenan en el medio producen, a su vez, electrones secundarios de menor energía que contribuyen a la dosis absorbida. Además, los fotones del haz interactúan con el cabezal del acelerador (colimador primario, filtro aplanador, mandíbulas, multiláminas) y con el propio maniquí, generando radiación dispersa con un espectro de energía más blando que el haz primario. El mecanismo espectral subyacente fue analizado en detalle por Benmakhlouf y Andreo~\cite{benmakhlouf2017}, quienes calcularon con el código PENELOPE la fluencia diferencial en energía de fotones y electrones dentro del volumen sensible de diversos detectores de campos pequeños para haces de 6~MV de un acelerador lineal. Sus resultados muestran que, al reducir el tamaño de campo, la contribución de fotones dispersados de baja energía provenientes del cabezal y del maniquí disminuye, produciendo un endurecimiento del espectro fotónico en el eje central. Este endurecimiento tiene una consecuencia directa sobre la distribución de energía de los electrones secundarios: en campos pequeños, la energía media de los electrones en el eje central es más alta que en campos grandes, porque hay menos fotones de baja energía que generen electrones secundarios lentos~\cite{benmakhlouf2017,palmans2018}.

La conexión con el \textit{quenching} es inmediata. Si la fracción de electrones con energía inferior a 125~keV cambia con el tamaño de campo, entonces la magnitud del \textit{quenching} también cambia con el tamaño de campo, y la respuesta del centelleador deja de ser independiente de las condiciones de irradiación. Valdes Santurio y Andersen~\cite{santurio2019} cuantificaron este efecto mediante el factor de corrección por \textit{quenching} de ionización, $k_{\mathrm{ioq}}$, definido como el cociente entre la luz que produciría un centelleador ideal (sin \textit{quenching}) y la luz que produce el centelleador real, ambas calculadas por Monte Carlo. Para un haz de 6~MV con el centelleador BCF-60, el $k_{\mathrm{ioq}}$ varía $(0{.}6 \pm 0{.}2)$\% entre campos de $0{.}5 \times 0{.}5$~cm$^{2}$ y $10 \times 10$~cm$^{2}$. El efecto es pequeño, pero estadísticamente significativo y sistemático: el $k_{\mathrm{ioq}}$ crece con el tamaño de campo porque los campos más grandes aportan más radiación dispersa de baja energía, lo que incrementa la fracción de electrones secundarios sujetos a \textit{quenching}~\cite{santurio2019}.

La dependencia con la calidad del haz es aún más pronunciada. Para un campo fijo de $10 \times 10$~cm$^{2}$, el $k_{\mathrm{ioq}}$ varía aproximadamente $(2 \pm 0{.}4)$\% entre 4~MV y 15~MV~\cite{santurio2019}. Esta variación es mayor que la dependencia con el tamaño de campo porque el cambio en la calidad del haz (medido por el TPR$_{20,10}$) modifica de forma más drástica la distribución de energía de los electrones secundarios: un haz de 4~MV tiene un espectro electrónico significativamente más blando que uno de 15~MV, con una fracción mayor de electrones por debajo de 125~keV~\cite{santurio2019}. La implicación práctica es que el \textit{quenching} tiene un impacto mayor cuando se utiliza un PSD para medir factores de corrección de calidad de haz ($k_{Q}$) que cuando se miden factores de salida de campo.

Estos resultados redefinen el papel del centelleador plástico en la dosimetría clínica. La variación de 0.6\% en los factores de salida de campo está dentro de la incertidumbre de las mediciones clínicas típicas, lo que explica por qué el efecto no había sido observado experimentalmente: estaba enmascarado por el ruido de medición. Sin embargo, la variación de $\sim$2\% en los factores $k_{Q}$ no es despreciable. Si un PSD se utiliza como instrumento de transferencia para verificar los factores $k_{Q}$ de cámaras de ionización --- una aplicación propuesta en la literatura ---, ignorar el \textit{quenching} introduciría un error sistemático del orden del 1--2\%, suficiente para comprometer la trazabilidad metrológica del resultado~\cite{santurio2019}.

El trabajo de Boivin \textit{et al.}~\cite{boivin2016} complementa esta perspectiva desde el extremo de baja energía. Utilizando tres modelos de centelleador (BCF-10, BCF-12 y BCF-60) expuestos a 13 calidades de haz entre 20~kVp y $^{60}$Co, los autores demostraron que la sensibilidad del PSD normalizada a la respuesta en $^{60}$Co disminuye de forma apreciable por debajo de 100~keV de energía media, con desviaciones de hasta un 5\% respecto al valor de referencia. La aplicación del modelo de Birks a los datos experimentales, integrando sobre el espectro de electrones secundarios, mostró un acuerdo razonable pero imperfecto: los valores de $kB$ extraídos dependen del rango de energía cubierto por el ajuste, sugiriendo que la asunción de $kB$ constante puede ser una aproximación limitada para haces con espectros muy diferentes al de calibración~\cite{boivin2016}.

%La imagen que emerge de estos trabajos es que el espectro de electrones secundarios --- una cantidad que raramente se examina en la práctica clínica cotidiana --- contiene la clave para entender el \textit{quenching} en haces de megavoltaje. La fracción de dosis depositada por electrones de baja energía no es una propiedad fija del haz, sino una variable que cambia con el tamaño de campo, con la calidad del haz, con la profundidad en el maniquí y con la presencia o ausencia del filtro aplanador. Cada una de estas variables modula el $k_{\mathrm{ioq}}$ de forma predecible mediante simulación Monte Carlo, pero invisible para una medición experimental que no separe la señal de centelleo de sus factores de influencia. % La sección siguiente muestra cómo el hallazgo de Gingras \textit{et al.} reveló que, en un detector específico, esta dependencia queda enmascarada por una compensación accidental.
% ─────────────────────────────────────────
\section{El hallazgo de Gingras: una compensación accidental}
% PRB-0002 en TrueBeam 6 MV. Separación Pscint y kioq.
% Compensación: factor total 0.999-1.002.
% La pregunta que deja abierta.
% Refs: gingras2025, papaconstadopoulos2014, casar2019

% CONTENIDO AQUÍ


El trabajo de Gingras \textit{et al.}~\cite{gingras2025} representa la primera determinación completa de los factores de corrección de salida de campo de un detector de centelleo plástico comercial, el PRB-0002 de la plataforma Hyperscint RP-200 de Medscint, utilizando el formalismo del TRS-483~\cite{palmans2018} adaptado a la dosimetría con centelleadores. La importancia de este trabajo no radica solamente en los valores numéricos obtenidos, sino en la separación explícita de cada factor de influencia que contribuye al factor de corrección total, y en lo que esa separación revela sobre la física del detector.

El marco formalista adoptado por Gingras \textit{et al.} se basa en la cadena de perturbación propuesta originalmente por Papaconstadopoulos \textit{et al.}~\cite{papaconstadopoulos2014} para detectores de centelleo y de estado sólido. La dosis absorbida en un punto de agua, $D^{f}_{w,Q}$, se relaciona con la lectura del detector a través de una cadena de cuatro factores de corrección:

\begin{equation}
D^{f}_{w,Q} = (k_{\mathrm{vol}})^{f}_{Q} \cdot (P_{\mathrm{scint}})^{f}_{Q} \cdot (P_{\mathrm{wall}})^{f}_{Q} \cdot D^{f}_{\mathrm{det},Q}
\label{eq:perturbation_chain}
\end{equation}

\noindent donde $k_{\mathrm{vol}}$ es el factor de corrección por promediado volumétrico (que da cuenta de la diferencia entre la dosis promediada en el volumen finito del detector y la dosis en un punto), $P_{\mathrm{scint}}$ es el factor de perturbación por composición del material centelleador (el cociente entre la dosis en un volumen de agua del mismo tamaño y la dosis en el volumen del centelleador real), $P_{\mathrm{wall}}$ es el factor de perturbación por la pared protectora, la fibra óptica y demás materiales no centelleadores del ensamblaje, y $D^{f}_{\mathrm{det},Q}$ es la dosis absorbida en el volumen sensible del detector completamente ensamblado~\cite{gingras2025}. Para un PSD, la cadena se extiende con un factor adicional que no tiene análogo en detectores convencionales: el factor de corrección por \textit{quenching} de ionización, $k_{\mathrm{ioq}}$, clasificado como un factor de influencia de tipo~II (dependiente de la calidad del haz). El factor de corrección de salida de campo del detector resulta entonces:

\begin{equation}
k^{f_{\mathrm{clin}},f_{\mathrm{msr}}}_{Q_{\mathrm{clin}},Q_{\mathrm{msr}}} = \left(k_{\mathrm{vol}} \cdot P_{\mathrm{scint}} \cdot P_{\mathrm{wall}} \cdot k_{\mathrm{ioq}}\right)^{f_{\mathrm{clin}},f_{\mathrm{msr}}}_{Q_{\mathrm{clin}},Q_{\mathrm{msr}}}
\label{eq:kfclin}
\end{equation}

\noindent donde la notación $(X)^{f_{\mathrm{clin}},f_{\mathrm{msr}}}_{Q_{\mathrm{clin}},Q_{\mathrm{msr}}}$ indica el cociente del factor $X$ evaluado en el campo clínico y en el campo de referencia específico de la máquina ($f_{\mathrm{msr}} = 10 \times 10$~cm$^{2}$ en el caso estudiado)~\cite{gingras2025}.

Gingras \textit{et al.} calcularon cada uno de estos factores por separado mediante simulaciones Monte Carlo con EGSnrc, utilizando los \textit{phase space files} del TrueBeam 6~MV de Varian y una geometría detallada del detector PRB-0002. El $k_{\mathrm{ioq}}$ se obtuvo con el \textit{ausgab object} \texttt{egs\_light\_scoring}, que implementa el modelo de Birks modificado de Valdes Santurio~\cite{santurio2020} con $kB = 0{.}019$~cm\,MeV$^{-1}$ y poder de frenado restringido con umbral $\Delta = 1$~keV. El $P_{\mathrm{scint}}$ se calculó como el cociente de dosis entre un volumen de agua y un volumen del centelleador real, ambos con la misma geometría, eliminando así la contribución de los materiales circundantes.

Los resultados para campos entre $0{.}6 \times 0{.}6$ y $30 \times 30$~cm$^{2}$ revelan un patrón notable. El $P_{\mathrm{scint}}$ del PRB-0002 decrece con el tamaño de campo: al ampliar el campo, la contribución de fotones dispersados de baja energía aumenta, y dado que el centelleador tiene un $Z_{\mathrm{eff}}$ ligeramente diferente al del agua, su interacción fotoeléctrica difiere levemente, produciendo una dosis relativa que se aparta de la unidad. El $k_{\mathrm{ioq}}$, por el contrario, crece con el tamaño de campo: los fotones dispersados de baja energía generan electrones secundarios más lentos que experimentan mayor \textit{quenching}, requiriendo una corrección creciente~\cite{gingras2025}. El resultado crucial es que estas dos tendencias opuestas se cancelan casi exactamente: el factor de corrección total $k^{f_{\mathrm{clin}},f_{\mathrm{msr}}}_{Q_{\mathrm{clin}},Q_{\mathrm{msr}}}$ del PRB-0002 se mantiene entre 0.999 y 1.002 en todo el rango de tamaños de campo, con una incertidumbre combinada del 0.5\%~\cite{gingras2025}. El $P_{\mathrm{wall}}$ resulta compatible con la unidad para todos los campos, y el $k_{\mathrm{vol}}$ es el único factor que contribuye de forma apreciable al apartarse de 1.

Este resultado tiene dos lecturas complementarias. La primera es práctica: el PRB-0002 se comporta como un detector casi ideal para la dosimetría de campos pequeños, con factores de corrección indistinguibles de la unidad dentro de la incertidumbre experimental. Esto valida la prescripción del TRS-483 para centelleadores plásticos, pero por razones diferentes a las asumidas originalmente. El TRS-483 asumía que el factor de corrección era unitario porque el detector era agua-equivalente; Gingras \textit{et al.} demuestran que el factor es unitario porque dos imperfecciones --- la perturbación material ($P_{\mathrm{scint}} \neq 1$) y el \textit{quenching} ($k_{\mathrm{ioq}} \neq 1$) --- se compensan mutuamente~\cite{gingras2025}. Este resultado tiene un antecedente directo en la literatura: Casar et al.~\cite{casar2019}
utilizaron el detector W1 PSD como referencia para la determinación de factores de salida
de campo en campos pequeños, asumiendo implícitamente un factor de corrección unitario.
El trabajo de Gingras et al.\ proporciona la justificación física de esa decisión práctica:
la prescripción no era incorrecta, sino que descansaba sobre una compensación cuyo mecanismo
no había sido identificado en 2019. La implicación operativa es que el PRB-0002 puede utilizarse como detector de referencia para determinar experimentalmente los factores de corrección de otros detectores con gran exactitud; los autores muestran que, usando este procedimiento, se pueden alcanzar incertidumbres del 1.9\% para un campo de $0{.}6 \times 0{.}6$~cm$^{2}$ y del 0.55\% para un campo de $2 \times 2$~cm$^{2}$~\cite{gingras2025}.

La segunda lectura es más profunda y constituye la motivación central del proyecto MDSIM1. La compensación entre $P_{\mathrm{scint}}$ y $k_{\mathrm{ioq}}$ observada en el PRB-0002 no es un resultado inevitable de la física de los centelleadores plásticos, sino una coincidencia numérica que depende de la composición específica de ese material centelleador en ese régimen de energía. El propio Gingras \textit{et al.} lo señalan al indicar que el \textit{quenching} de ionización tiene un impacto medible en haces de megavoltaje y que la compensación que observan es específica del PRB-0002~\cite{gingras2025}. $P_{\mathrm{scint}}$ depende del número atómico efectivo y la densidad del centelleador, que determinan su sección eficaz de interacción fotoeléctrica y de dispersión relativa al agua. $k_{\mathrm{ioq}}$ depende del parámetro de Birks $kB$, que es una propiedad de la composición molecular del polímero y del dopante, relacionada en última instancia con el radio de Förster del material. Los dos factores están gobernados por mecanismos físicos distintos con dependencias funcionales independientes de los parámetros del material.

Esta independencia implica que un cambio en la composición del centelleador --- por ejemplo, un polímero base diferente, una concentración de dopante modificada o un aditivo para mejorar la resistencia a la radiación --- alterará $P_{\mathrm{scint}}$ y $k_{\mathrm{ioq}}$ de forma diferente, y la compensación observada en el PRB-0002 puede romperse. Si esto ocurre, el nuevo material tendría un factor de corrección significativamente diferente de la unidad, y su uso clínico requeriría correcciones explícitas. Si, por el contrario, se diseña un material que mantiene la compensación, se obtiene un detector intrínsecamente auto-corregido.

El antecedente más cercano al enfoque de diseño compensado es el trabajo de Papaconstadopoulos \textit{et al.}~\cite{papaconstadopoulos2014}, quienes mostraron mediante simulaciones Monte Carlo con EGSnrc que, para algunos detectores de estado sólido (diodos, diamante), la modificación de la geometría del volumen sensible puede compensar la perturbación material con el efecto de promediado volumétrico, produciendo factores de corrección cercanos a la unidad para campos específicos. La diferencia conceptual con lo observado por Gingras \textit{et al.} es que en el caso del PSD la compensación no involucra el promediado volumétrico sino un efecto puramente material: la perturbación por composición ($P_{\mathrm{scint}}$) y la no-linealidad de la señal por \textit{quenching} ($k_{\mathrm{ioq}}$). Esta compensación material-\textit{quenching} no había sido reportada previamente y abre una línea de investigación nueva: la posibilidad de optimizar la composición del centelleador para que el producto $P_{\mathrm{scint}} \cdot k_{\mathrm{ioq}}$ sea lo más cercano posible a la unidad en un rango amplio de condiciones clínicas.

La pregunta que el hallazgo de Gingras deja abierta puede formularse de manera precisa: ¿para qué materiales centelleadores, en qué régimen energético y para qué geometrías de detector se cumple que $P_{\mathrm{scint}}^{f_{\mathrm{clin}},f_{\mathrm{msr}}} \cdot k_{\mathrm{ioq}}^{f_{\mathrm{clin}},f_{\mathrm{msr}}} \approx 1$? Responder esta pregunta requiere una herramienta que permita variar la composición del material de forma sistemática y calcular ambos factores separadamente. Esa herramienta es MDSIM1.

\section{La contribución de MDSIM1}
% Geant4 + phase space TrueBeam + modos Birks/Optical.
% Métrica Cmaterial. No es replicación, es exploración.
% Refs: mdsim2026, dietzlaursonn2016, poon2005

% CONTENIDO AQUÍ

%MDSIM1 (\textit{MultiDetector Simulation 1}) es una plataforma de simulación Monte Carlo construida en Geant4 que ha sido diseñada específicamente para calcular de forma separada los factores $P_{\mathrm{scint}}$ y $k_{\mathrm{ioq}}$ de detectores de centelleo plástico en haces clínicos de fotones y electrones. Su diseño responde directamente a las cuatro brechas identificadas en la sección anterior, y su arquitectura refleja las decisiones metodológicas necesarias para abordarlas con rigor.

%La fuente de radiación de MDSIM1 es un archivo de espacio de fases (\textit{phase space}) del acelerador Varian TrueBeam de 6~MV, descargado del repositorio oficial de Varian y validado contra datos experimentales de curvas de dosis en profundidad y perfiles laterales antes de incorporar cualquier componente de centelleo. Los archivos de espacio de fases se leen en formato IAEA y se propagan a través de un sistema de colimación simulado con las dimensiones nominales del cabezal, lo que permite generar campos cuadrados desde $0{.}5 \times 0{.}5$ hasta $30 \times 30$~cm$^{2}$. Esta configuración reproduce las condiciones exactas utilizadas por Valdes Santurio y Andersen~\cite{santurio2019} y por Gingras \textit{et al.}~\cite{gingras2025} en sus cálculos con EGSnrc, lo que hace posible una comparación directa de resultados.

%El maniquí de irradiación es un volumen de agua paralelepipédico con dimensiones suficientes para garantizar retrodispersión completa. Los detectores se posicionan en su interior a la profundidad de referencia ($d_{\mathrm{ref}} = 10$~cm a SSD~$= 90$~cm, o $d_{\mathrm{max}}$ según la aplicación). MDSIM1 implementa una arquitectura modular de detectores que permite insertar y configurar distintas geometrías de volumen sensible sin modificar el código de transporte. Los módulos de detector actualmente implementados incluyen geometrías básicas parametrizables --- cubo, cilindro y esfera, con material, dimensiones y posición configurables en tiempo de ejecución --- y el módulo \texttt{model11}, que importa la geometría completa del detector Blue Physics Model~11 desde un archivo GDML (\textit{Geometry Description Markup Language}) generado a partir del modelo CAD del fabricante~\cite{bauer2026}. El módulo \texttt{model11} permite seleccionar por nombre los volúmenes del ensamble que actúan como volúmenes sensibles, posicionar el detector en la superficie o en profundidad, y activar una función de corte en la interfaz aire-agua que reparte automáticamente la geometría del detector entre los dos medios cuando el detector cruza la superficie del maniquí.

%La lista de física utilizada en MDSIM1 para fotones MV es \texttt{G4EmStandardPhysics\_option4}, que incluye los modelos de dispersión múltiple \texttt{G4WentzelVIModel}, ionización con correcciones de subcapas, y producción de pares con efecto LPM. Esta lista ha sido validada extensamente para aplicaciones de dosimetría en radioterapia y ofrece el mejor compromiso entre precisión y costo computacional para el régimen de energías relevante~\cite{poon2005}. Para el modo Optical, la lista se extiende con \texttt{G4OpticalPhysics}, que habilita la generación de fotones de centelleo (\texttt{G4Scintillation}), la absorción óptica (\texttt{G4OpAbsorption}), los procesos de interfaz óptica (\texttt{G4OpBoundaryProcess}) y la generación de radiación Cherenkov (\texttt{G4Cerenkov}).

%MDSIM1 opera en dos modos complementarios que abordan aspectos diferentes de la física del detector:

%El \textbf{modo Birks} calcula $k_{\mathrm{ioq}}$ a partir de la deposición de energía en el volumen sensible, sin simular el transporte de fotones ópticos. En cada paso de simulación de un electrón dentro del centelleador, la clase \texttt{G4EmSaturation} aplica el factor de Birks utilizando el $\mathrm{d}E/\mathrm{d}x$ local del paso y el valor de $kB$ definido por el usuario. La implementación utiliza dos acumuladores simultáneos en cada evento: uno registra la energía depositada total $E_{\mathrm{dep}}$ (proporcional a $L_{\mathrm{ideal}}$, la luz que se produciría sin \textit{quenching}) y otro registra la energía depositada ponderada por el factor de Birks paso a paso (proporcional a $L_{\mathrm{quenched}}$, la luz real). El factor de corrección por \textit{quenching} se obtiene como $k_{\mathrm{ioq}} = \sum E_{\mathrm{dep},i} / \sum (E_{\mathrm{dep},i} \cdot f_{\mathrm{Birks},i})$, donde la suma recorre todos los pasos de todos los eventos. Esta arquitectura de doble acumulador correlacionado es estadísticamente más eficiente que la alternativa de dos corridas independientes (una con $kB = 0$ y otra con $kB$ real), porque el cociente de sumas correlacionadas tiene una varianza significativamente menor que el cociente de medias independientes~\cite{ebenau2016}. El $P_{\mathrm{scint}}$ se calcula por separado como el cociente de dosis absorbida entre un volumen de agua y un volumen del centelleador con la misma geometría.

%El \textbf{modo Optical} simula el transporte completo de los fotones de centelleo desde su generación en el volumen sensible hasta su detección en el extremo de la fibra óptica. Este modo es necesario para aplicaciones donde la señal Cherenkov es relevante --- como en haces de electrones, donde Bauer \textit{et al.}~\cite{bauer2026} reportan que la contribución Cherenkov puede alcanzar el 50\% de la señal total en un campo de $10 \times 10$~cm$^{2}$ --- y para estudios donde la eficiencia de colección óptica del detector influye en la señal medida. El modo Optical requiere la definición de propiedades ópticas del centelleador (índice de refracción, longitud de absorción, espectro de emisión, rendimiento de centelleo y constante de decaimiento) como función de la energía del fotón óptico.

La contribución científica de MDSIM1 no es replicar lo que Gingras \textit{et al.}~\cite{gingras2025} ya han hecho con EGSnrc, sino extender su hallazgo puntual a un marco generalizable. El concepto central que articula esta extensión es la métrica de compensación material, $C_{\mathrm{material}}$, definida como:

\begin{equation}
C_{\mathrm{material}} = P_{\mathrm{scint}}^{f_{\mathrm{clin}},f_{\mathrm{ref}}} \cdot k_{\mathrm{ioq}}^{f_{\mathrm{clin}},f_{\mathrm{ref}}}
\label{eq:cmaterial}
\end{equation}

Gingras \textit{et al.} demostraron que $C_{\mathrm{material}} \approx 1$ para el PRB-0002 en 6~MV. MDSIM1 permite evaluar $C_{\mathrm{material}}$ para cualquier material centelleador definido en Geant4, en cualquier régimen energético accesible mediante un archivo de espacio de fases, y para cualquier geometría de detector importable en GDML. Un material con $C_{\mathrm{material}} \approx 1$ en un rango amplio de tamaños de campo y calidades de haz es un candidato intrínsecamente auto-corregido para dosimetría clínica; un material con $C_{\mathrm{material}}$ significativamente diferente de la unidad requiere correcciones explícitas que dependen de las condiciones de irradiación.

Para que los valores de $C_{\mathrm{material}}$ calculados por MDSIM1 sean creíbles, la plataforma debe superar dos pruebas de validación secuenciales. La primera es la reproducción de la dosis absorbida en agua: las curvas de dosis en profundidad y los perfiles laterales calculados con MDSIM1 deben concordar con los datos experimentales y con los resultados de EGSnrc dentro de las tolerancias establecidas por Poon y Verhaegen~\cite{poon2005} (menos del 2\% en la región de equilibrio). La segunda es la reproducción del $k_{\mathrm{ioq}}$ de referencia: el valor obtenido con MDSIM1 para BCF-60 en un campo de $10 \times 10$~cm$^{2}$ a 6~MV debe ser consistente con el $\sim$0.6\% reportado por Valdes Santurio y Andersen~\cite{santurio2019}, dentro de las incertidumbres combinadas de ambos cálculos y de la diferencia esperada entre esquemas de transporte. Solo después de superar estas dos pruebas puede MDSIM1 extenderse a materiales y regímenes no explorados previamente.

%Es importante delimitar lo que MDSIM1 no pretende resolver. No determina el $kB$ del material de Blue Physics experimentalmente: asume $kB = 0{.}019$~cm\,MeV$^{-1}$ como punto de partida y evalúa la sensibilidad del resultado frente a variaciones de este parámetro ($kB = 0{.}010$, 0.019 y 0.025~cm\,MeV$^{-1}$). No sustituye la medición experimental del factor de corrección; proporciona la herramienta de predicción que orienta el diseño experimental. Y no resuelve la limitación fundamental del modelo de Birks --- su naturaleza semiempírica y material-dependiente --- sino que la utiliza como el modelo de mejor compromiso entre aplicabilidad práctica y solidez teórica para el régimen de interés.

% ─────────────────────────────────────────

\newpage

% ─────────────────────────────────────────

\section{El espacio de materiales como ventaja competitiva}
% Cmaterial ≈ 1 como criterio de diseño.
% Bauer 2026: Model 11. Blue Physics.
% Conexión Ian Rossi: detector portátil estado sólido.
% Refs: bauer2026, poeschl2020, laplace2022

% CONTENIDO AQUÍ



La métrica $C_{\mathrm{material}}$ definida en la ecuación~\eqref{eq:cmaterial} transforma una observación experimental puntual --- la compensación entre $P_{\mathrm{scint}}$ y $k_{\mathrm{ioq}}$ en el PRB-0002 --- en un criterio de diseño cuantificable. Un material centelleador con $C_{\mathrm{material}} \approx 1$ en un rango amplio de tamaños de campo y calidades de haz es intrínsecamente auto-corregido: puede utilizarse en la clínica sin aplicar factores de corrección detector-específicos, reduciendo la cadena de incertidumbre de la medición y simplificando el procedimiento dosimétrico. Un material con $C_{\mathrm{material}}$ significativamente diferente de la unidad requiere correcciones que dependen de las condiciones de irradiación, lo que limita su aplicabilidad práctica y aumenta el riesgo de error sistemático.

Esta perspectiva convierte la simulación Monte Carlo en una herramienta de prospección de materiales. En lugar de fabricar y caracterizar experimentalmente cada formulación candidata --- un proceso costoso, lento y que requiere acceso a instalaciones de irradiación calibradas ---, MDSIM1 puede evaluar $C_{\mathrm{material}}$ para cualquier centelleador cuya composición elemental, densidad y $kB$ sean conocidos o estimables. El espacio de parámetros relevante incluye el polímero base (poliestireno, polivinilueno, PMMA, polisulfona), el tipo y la concentración de dopante fluorescente (p-terfenilo, POPOP, derivados de oxazol), la presencia de aditivos funcionales (estabilizadores UV, agentes endurecedores, antioxidantes) y la densidad final del compuesto. Cada combinación define un punto en el espacio de materiales, con un $Z_{\mathrm{eff}}$, una densidad y un $kB$ que determinan $P_{\mathrm{scint}}$ y $k_{\mathrm{ioq}}$ respectivamente.

El contexto industrial que motiva esta exploración es la empresa Blue Physics, que desarrolla detectores de centelleo plástico de nueva generación con materiales propios. El detector Blue Physics Model~11 ha sido caracterizado experimentalmente por Bauer \textit{et al.}~\cite{bauer2026} en haces de electrones de 6--12~MeV en un acelerador Elekta Synergy. Los resultados son alentadores: el detector demuestra linealidad (coeficiente de variación de 1.3\% entre 1 y 1000~UM), independencia de tasa de dosis (1.0\%), repetibilidad (0.3\%), isotropía (0.8\%) y un factor de calibración estable de $(2{.}29 \pm 0{.}01)$~cGy/$\mu$C~\cite{bauer2026}. Las curvas de dosis en profundidad medidas con el Model~11 concuerdan con las medidas con cámaras de ionización de referencia dentro de 1~mm en $R_{80}$ y $R_{50}$, y los factores de salida de campo para campos de $3 \times 3$ a $7 \times 7$~cm$^{2}$ presentan una desviación media de $(0{.}4 \pm 0{.}8)$\% respecto a la cámara Advanced Markus~\cite{bauer2026}. Las mediciones de dosis en superficie en un maniquí RW3 y en un maniquí antropomórfico muestran un error absoluto medio inferior al 1.5\% y al 1.4\%, respectivamente, comparado con cámaras de ionización y películas radiocrómicas~\cite{bauer2026}.

Estos resultados confirman que el detector Blue Physics tiene un desempeño clínico competitivo. Sin embargo, la caracterización de Bauer \textit{et al.} fue realizada exclusivamente en haces de electrones, un régimen donde la contribución Cherenkov en la fibra óptica es sustancial (aproximadamente el 50\% de la señal total en un campo de $10 \times 10$~cm$^{2}$, corregida mediante un cociente de canal adyacente ACR~$= 0{.}916 \pm 0{.}013$)~\cite{bauer2026}. En haces de fotones MV, donde el \textit{quenching} de ionización es el factor de influencia dominante, el detector no ha sido caracterizado de la misma manera que el PRB-0002 por Gingras \textit{et al.}~\cite{gingras2025}. Esto significa que la pregunta central --- si el material centelleador del Model~11 exhibe la misma compensación $P_{\mathrm{scint}} \cdot k_{\mathrm{ioq}} \approx 1$ que el PRB-0002 en haces de fotones MV --- permanece abierta. MDSIM1 fue diseñado precisamente para responderla.

La geometría del Model~11 ya está implementada en MDSIM1 como un módulo GDML importado directamente del modelo CAD del fabricante. El volumen sensible es un cilindro de 1~mm de diámetro y 1~mm de longitud (0.785~mm$^{3}$), confirmado por Bauer \textit{et al.}~\cite{bauer2026}. La simulación permite posicionar el detector en profundidad o en la superficie del maniquí de agua, con una función de corte automático en la interfaz aire-agua que reparte la geometría del detector entre ambos medios. Esta funcionalidad es relevante para la dosimetría en superficie y la dosimetría \textit{in vivo}, dos aplicaciones que Bauer \textit{et al.} identifican como las más prometedoras del Model~11.

El valor de $kB$ del material centelleador de Blue Physics no ha sido determinado experimentalmente. Toda la simulación parte del valor de referencia $kB = 0{.}019$~cm\,MeV$^{-1}$ medido para BCF-60~\cite{santurio2019}, con un análisis de sensibilidad que evalúa la variación de $k_{\mathrm{ioq}}$ y de $C_{\mathrm{material}}$ frente a $kB = 0{.}010$ y $kB = 0{.}025$~cm\,MeV$^{-1}$. Si la simulación muestra que $C_{\mathrm{material}}$ es sensible a $kB$ en el rango explorado --- como sugieren los resultados de Pöschl \textit{et al.}~\cite{poeschl2020} para SCSF-78 y BC-408 ---, entonces la determinación experimental del $kB$ propio del material Blue Physics se convierte en una prioridad. La metodología modelo-independiente de Pöschl, basada en un detector segmentado expuesto a protones de 30--100~MeV, ofrece un camino para esa medición, aunque requiere acceso a una instalación de protonterapia.

%La conexión con el proyecto de Ian Rossi emerge naturalmente de esta perspectiva. El desarrollo de un detector portátil de radiación basado en estado sólido requiere seleccionar materiales centelleadores que optimicen simultáneamente la eficiencia de detección, la robustez mecánica y la estabilidad temporal. Si el \textit{quenching} no se considera en el diseño del material, un detector portátil podría exhibir una respuesta no lineal con la energía de la radiación incidente que comprometería su exactitud. El modelo FRET de Ogawa \textit{et al.}~\cite{ogawa2018} sugiere que el radio de Förster $R_f$ --- que depende de la arquitectura molecular del centelleador --- es el parámetro fundamental que controla la susceptibilidad al \textit{quenching}. Un material con $R_f$ pequeño (dopantes con menor eficiencia de transferencia dipolo-dipolo) será menos susceptible al \textit{quenching} pero posiblemente menos eficiente en la producción de luz; un material con $R_f$ grande será más eficiente pero más vulnerable. La exploración computacional de este compromiso --- que MDSIM1 hace posible a través de la variación paramétrica de $kB$ para distintas composiciones --- proporciona criterios de selección informados antes de comprometer recursos en la síntesis y caracterización de materiales.

Por tanto, el espacio de materiales centelleadores no es un inventario estático de compuestos comerciales con propiedades fijas, sino un paisaje paramétrico donde $Z_{\mathrm{eff}}$, densidad, $kB$ y $R_f$ interactúan para determinar la calidad dosimétrica del detector. %La empresa que comprenda ese paisaje y sepa navegarlo con herramientas predictivas tiene una ventaja competitiva sobre la que opera por ensayo y error. MDSIM1 es la primera herramienta diseñada específicamente para mapear ese paisaje en el contexto de la dosimetría con centelleadores plásticos.

% ─────────────────────────────────────────

\section{Conclusiones}

\subsection*{Estado actual y preguntas abiertas}
% Checklist MDSIM1. Qué falta. Próximos pasos.
% Cierre con la pregunta más importante, no con conclusión.

% CONTENIDO AQUÍ


La revision del estado del arte establece un mapa claro: el \textit{quenching} de ionización en centelleadores plásticos es un fenómeno cuantificable, modelable y --- en al menos un caso específico --- accidentalmente compensado por la perturbación material del detector. Las herramientas teóricas existen: la ley de Birks y su extensión modificada por Valdes Santurio~\cite{santurio2020}, los modelos moleculares basados en FRET~\cite{ogawa2018}, y los códigos Monte Carlo capaces de implementar ambos. Los datos de referencia también existen: las mediciones y simulaciones de Gingras \textit{et al.}~\cite{gingras2025} para el PRB-0002 en 6~MV, las caracterizaciones espectrales de Valdes Santurio y Andersen~\cite{santurio2019} para BCF-60, y la validación experimental de Bauer \textit{et al.}~\cite{bauer2026} para el Blue Physics Model~11 en haces de electrones. Lo que falta es la conexión entre ellos: una plataforma que use esas herramientas teóricas, se valide contra esos datos de referencia, y extienda los resultados al territorio inexplorado de otros materiales, otros regímenes energéticos y otro código Monte Carlo.

%A la fecha, la infraestructura de MDSIM1 se encuentra en un estado de desarrollo avanzado. Los componentes completados incluyen la simulación funcional del haz a partir de los archivos de espacio de fases del TrueBeam 6~MV, la implementación del maniquí de agua con dimensiones de referencia, los módulos de detector básicos (cubo, cilindro y esfera) con material, tamaño y posición configurables en tiempo de ejecución, el módulo \texttt{model11} con geometría importada desde GDML, el sistema de \textit{scoring} de dosis absorbida, y los procedimientos de calibración y barridos de curvas de dosis en profundidad en modo \textit{batch}. Los componentes en desarrollo activo incluyen la implementación del modo Birks (\texttt{G4EmSaturation} con doble acumulador correlacionado $L_{\mathrm{ideal}}/L_{\mathrm{quenched}}$), la validación de convergencia del cociente frente a variaciones del parámetro de control de paso \texttt{dRoverRange}, y la comparación cuantitativa de $k_{\mathrm{ioq}}$ entre Geant4 y los valores de referencia de EGSnrc.

MDSIM1 tiene un camino a seguir: Primero, verificar que la simulación con $kB = 0$ produce $k_{\mathrm{ioq}} = 0$ exactamente --- una prueba de consistencia interna que valida la implementación del doble acumulador. Segundo, verificar que la simulación con $kB = 0{.}019$~cm\,MeV$^{-1}$ en un campo de $10 \times 10$~cm$^{2}$ a 6~MV reproduce un $k_{\mathrm{ioq}} \approx 0{.}6$\%, consistente con el valor reportado por Valdes Santurio y Andersen~\cite{santurio2019}. Tercero, completar el barrido de \texttt{dRoverRange} (0.1, 0.05, 0.02, 0.01) y demostrar convergencia del cociente $L_{\mathrm{quenched}}/L_{\mathrm{ideal}}$ --- no de cada término por separado, sino del cociente, que es la cantidad físicamente relevante~\cite{ebenau2016}. Cuarto, calcular $P_{\mathrm{scint}}$ y $k_{\mathrm{ioq}}$ separadamente para el material del PRB-0002 en el rango de campos de $0{.}6 \times 0{.}6$ a $30 \times 30$~cm$^{2}$ y comparar el $C_{\mathrm{material}}$ resultante con los valores de Gingras \textit{et al.}~\cite{gingras2025}. Quinto, extender el cálculo a composiciones de material centelleador alternativas, evaluando $C_{\mathrm{material}}$ para al menos tres valores de $kB$ (0.010, 0.019 y 0.025~cm\,MeV$^{-1}$) que cubran el rango reportado en la literatura para centelleadores de base poliestireno y polivinilueno~\cite{boivin2016,poeschl2020}.

Este estado del arte no concluye con una afirmación, sino con la pregunta que lo motivó desde el inicio y que permanece abierta: si el \textit{quenching} de ionización no es un defecto a ignorar sino un parámetro a diseñar, ¿cuál es la composición del centelleador que minimiza la necesidad de corrección en el rango más amplio posible de condiciones clínicas? La respuesta (hasta la fecha) no está en ningún paper publicado. Potencialmente está en la simulación que MDSIM1 se propone completar.

\subsection*{Lo que el estado del arte no responde}
% Tabla de 4 brechas: Geant4 vs EGSnrc, universalidad,
% régimen kV, control de paso.
% Refs: poon2005, dietzlaursonn2016, boivin2016, gingras2025

% CONTENIDO AQUÍ


%Las secciones anteriores han trazado una línea narrativa que va desde el descubrimiento del \textit{quenching} de ionización en centelleadores orgánicos hasta su cuantificación en haces clínicos de megavoltaje y la demostración de una compensación accidental entre $P_{\mathrm{scint}}$ y $k_{\mathrm{ioq}}$ en un detector específico. Esta línea de investigación, desarrollada principalmente por dos grupos --- el de Andersen y Valdes Santurio en la DTU de Dinamarca (modelado del \textit{quenching}) y el de Archambault, Beaulieu y Gingras en la Université Laval y Medscint en Canadá (caracterización experimental y Monte Carlo del PRB-0002) ---, ha sido extraordinariamente productiva. Sin embargo, deja abiertas cuatro preguntas que definen el espacio de contribución del proyecto MDSIM1.

%La Tabla~\ref{tab:brechas} resume estas cuatro brechas, su origen en la literatura y la forma en que MDSIM1 propone abordarlas.

\begin{table}[H]
\centering
\caption{Brechas identificadas en el estado del arte y posicionamiento de MDSIM1.}
\label{tab:brechas}
\small
\begin{tabular}{p{4.2cm} p{3.0cm} p{3.5cm} p{3.5cm}}
\toprule
\textbf{Pregunta abierta} & \textbf{Origen} & \textbf{Estado en la literatura} & \textbf{Contribución de MDSIM1} \\
\midrule
¿Geant4 reproduce los valores de $k_{\mathrm{ioq}}$ calculados con EGSnrc para el mismo detector y haz? &
Implícita en Gingras~\cite{gingras2025} y Poon y Verhaegen~\cite{poon2005} &
Sin datos: todos los cálculos de $k_{\mathrm{ioq}}$ para PSD en haces MV han sido realizados con EGSnrc &
Primer cálculo pareado de $k_{\mathrm{ioq}}$ con Geant4 usando \texttt{G4EmSaturation}, en las mismas condiciones de Gingras \\
\addlinespace
¿La compensación $P_{\mathrm{scint}} \cdot k_{\mathrm{ioq}} \approx 1$ es universal o depende de la composición del material? &
Gingras~\cite{gingras2025}, Pöschl \textit{et al.}~\cite{poeschl2020} &
Sin datos: solo se ha demostrado para el PRB-0002 (BCF-60). Pöschl muestra que $kB$ y la función de \textit{quenching} son material-dependientes &
Evaluación de $C_{\mathrm{material}} = P_{\mathrm{scint}}^{f,\mathrm{ref}} \cdot k_{\mathrm{ioq}}^{f,\mathrm{ref}}$ para múltiples composiciones de centelleador \\
\addlinespace
¿El $k_{\mathrm{ioq}}$ varía con el régimen energético (kV \textit{vs.} MV) de forma material-específica? &
Boivin \textit{et al.}~\cite{boivin2016}, Valdes Santurio~\cite{santurio2020} &
Parcial: Boivin muestra variación de sensibilidad en kV pero con $kB$ constante; Valdes Santurio valida en kV solo para BCF-60 &
Comparación kV \textit{vs.} MV con separación explícita $P_{\mathrm{scint}}/k_{\mathrm{ioq}}$ para materiales con distinto $kB$ \\
\addlinespace
¿El control de paso del transporte electrónico en Geant4 introduce un sesgo sistemático en el cálculo de $k_{\mathrm{ioq}}$? &
Dietz-Laursonn~\cite{ebenau2016} &
Documentado para partículas pesadas en contexto de física de altas energías. No cuantificado para fotones MV en PSD &
Análisis de convergencia de $k_{\mathrm{ioq}}$ \textit{vs.} \texttt{dRoverRange} en Geant4 \\
\bottomrule
\end{tabular}
\end{table}

%La primera brecha es de verificación cruzada entre códigos Monte Carlo. Toda la cadena de cálculo de $k_{\mathrm{ioq}}$ para centelleadores plásticos en haces de fotones MV ha sido desarrollada exclusivamente con EGSnrc: el \textit{ausgab object} \texttt{egs\_light\_scoring} de Valdes Santurio~\cite{santurio2020}, utilizado por Gingras \textit{et al.}~\cite{gingras2025}, implementa el modelo de Birks modificado con poder de frenado restringido ($\Delta = 1$~keV) y se beneficia del esquema de transporte electrónico optimizado de EGSnrc para cavidades. Geant4 implementa la corrección de Birks a través de la clase \texttt{G4EmSaturation}, que aplica el factor de saturación paso a paso utilizando el $\mathrm{d}E/\mathrm{d}x$ local de cada paso de simulación. Los dos enfoques son matemáticamente equivalentes en el límite de paso infinitesimal, pero operativamente diferentes porque dependen de esquemas de transporte distintos~\cite{ebenau2016}. Poon y Verhaegen~\cite{poon2005} establecieron que las diferencias entre Geant4 y EGSnrc en dosis absorbida en agua son menores al 2\% en la región de equilibrio de partículas cargadas, pero pueden ser significativamente mayores en la región de \textit{buildup}. Para $k_{\mathrm{ioq}}$, que es una cantidad derivada del cociente entre luz con y sin \textit{quenching}, no existe ninguna comparación publicada entre los dos códigos.

%La segunda brecha es de universalidad. La compensación entre $P_{\mathrm{scint}}$ y $k_{\mathrm{ioq}}$ reportada por Gingras \textit{et al.}~\cite{gingras2025} es un resultado para un material específico (el centelleador del PRB-0002) en un régimen específico (6~MV, campos de $0{.}6 \times 0{.}6$ a $30 \times 30$~cm$^{2}$). $P_{\mathrm{scint}}$ depende del número atómico efectivo y la densidad del centelleador; $k_{\mathrm{ioq}}$ depende del parámetro de Birks $kB$, que es una propiedad de la microestructura molecular del material. Los trabajos de Pöschl \textit{et al.}~\cite{poeschl2020} y Laplace \textit{et al.}~\cite{laplace2022} demuestran que la función de \textit{quenching} difiere sustancialmente entre materiales, incluso entre centelleadores comerciales basados en polímeros similares. No existe en la literatura un estudio que evalúe sistemáticamente si la compensación se mantiene, se refuerza o se destruye al variar la composición del centelleador. Esta pregunta no puede responderse experimentalmente sin fabricar y caracterizar múltiples prototipos, pero puede abordarse computacionalmente si se dispone de una plataforma de simulación que calcule $P_{\mathrm{scint}}$ y $k_{\mathrm{ioq}}$ separadamente para distintas definiciones de material.

%La tercera brecha es de régimen energético. Boivin \textit{et al.}~\cite{boivin2016} demostraron que la sensibilidad de los PSD varía apreciablemente en el régimen de kilovoltaje, con desviaciones de hasta un 5\% respecto al valor de referencia en $^{60}$Co para energías medias por debajo de 100~keV. Valdes Santurio y Andersen~\cite{santurio2020} validaron su modelo modificado de Birks en el régimen kV para BCF-60, pero con un único valor de $kB$. Dado que Boivin~\cite{boivin2016} reportó que la asunción de $kB$ constante puede ser limitada para haces con espectros muy diferentes al de calibración, la extensión al régimen kV con materiales distintos a BCF-60 requiere una evaluación separada. Para el contexto de Blue Physics, esta brecha es relevante porque el detector Model~11 ha sido caracterizado experimentalmente por Bauer \textit{et al.}~\cite{bauer2026} en haces de electrones de 6--12~MeV, pero su comportamiento en regímenes de kilovoltaje --- donde la fracción de dosis depositada por electrones de baja energía es mayor y el \textit{quenching} es más pronunciado --- no ha sido evaluado.

%La cuarta brecha es de implementación numérica. Dietz-Laursonn~\cite{ebenau2016} documentó que el $kB$ efectivo recuperado de una simulación Geant4 depende del tamaño de paso del transporte electrónico, como consecuencia directa de la no-linealidad de la ley de Birks. Pasos grandes sobreestiman el $\mathrm{d}E/\mathrm{d}x$ promedio por paso y, por tanto, sobreestiman el \textit{quenching}; pasos finos capturan mejor las fluctuaciones locales pero incrementan el costo computacional. En Geant4, los parámetros relevantes son \texttt{dRoverRange} y \texttt{finalRange} del modelo de dispersión múltiple (\texttt{G4UrbanMscModel} o \texttt{G4WentzelVIModel}). Este efecto ha sido documentado en el contexto de detectores de partículas para física de altas energías, pero no ha sido cuantificado para el caso específico de fotones MV en detectores PSD clínicos. Sin un estudio de convergencia que demuestre que el cociente $L_{\mathrm{quenched}}/L_{\mathrm{ideal}}$ es estable frente a variaciones del control de paso, cualquier valor de $k_{\mathrm{ioq}}$ obtenido con Geant4 carece de una base numérica verificada.

%Estas cuatro brechas no son independientes: están conectadas por una estructura lógica que va de lo particular a lo general. La primera brecha (verificación Geant4 \textit{vs.} EGSnrc) es un prerrequisito metodológico: antes de explorar la dependencia material o energética del $k_{\mathrm{ioq}}$, es necesario demostrar que la herramienta de cálculo reproduce los resultados de referencia. La cuarta brecha (control de paso) es un prerrequisito numérico de la primera: la verificación cruzada solo es significativa si el resultado de Geant4 ha convergido. La segunda brecha (universalidad) es la pregunta científica central, y la tercera (régimen kV) es su extensión natural hacia el dominio donde el \textit{quenching} es más pronunciado y la compensación más frágil.

% ─────────────────────────────────────────


\newpage
\bibliographystyle{unsrtnat}
\bibliography{refs}
% ============================================================
% referencias.bib — Estado del arte MDSIM1
% ============================================================



\end{document}